\chapter{Інтегроване моделювання й числові методи}\label{ch:modelling}

\section{Ієрархія моделей}\label{sec:11-hierarchy}

\begin{outofcorpus}[розділення масштабів]
У токамаку з $\Bzero\sim\SI{5}{Тл}$, $T\sim\SI{10}{кеВ}$ характерні часи впорядковані так: $\Omc{e}^{-1}\sim10^{-12}$~с $\ll\Omc{i}^{-1}\sim10^{-9}$~с $\ll\tau_{\mathrm A}=a/\vA\sim10^{-7}$--$10^{-6}$~с $\ll$ час кореляції дрейфової турбулентності $\sim10^{-5}$--$10^{-4}$~с $\ll\tauE\sim1$~с $\ll\tau_R\gtrsim10^{3}$~с. Довжини: $\rho_e\sim10^{-5}$~м $\ll\rho_i\sim10^{-3}$~м $\ll a\sim1$~м. Усереднення прибирає швидкий масштаб: гірофаза~$\to$ гірокінетика; альфвенівська динаміка~$\to$ послідовність рівноваг GS; магнітна поверхня~$\to$ 1.5D; профіль~$\to$ 0D. Ціна~--- замикання ззовні (перенесення, закон Ома, граничні умови).
\end{outofcorpus}
Межу розділення видно на GTC: турбулентність має масштаб гірорадіуса, а її зональні потоки~--- масштаб плазми, тож потрібен глобальний код~\cite[с.~4]{Wang2006}. Коди корпусу~--- у табл.~\ref{tab:11-codes}.

\begin{table}[htbp]\centering\footnotesize
\caption{Коди корпусу за ієрархією моделей.}\label{tab:11-codes}
\begin{tabularx}{\textwidth}{@{}l X l l l@{}}
\toprule
Код & Модель & Вимір & Джерело & Розд. \\
\midrule
глобальна модель & баланс частинок і енергії тліючого розряду & 0D & \cite{Martsenyuk2025} & \ref{ch:wall}\\
STREAM (DREAM) & прогоряння, утікачі & 0D & \cite{Hoppe2022} & \ref{ch:disruptions}\\
TASK/TR & перенесення сортів, струм, нагрів пучком & 1D (1.5D) & \cite{Fukuyama1992} & \ref{ch:transport}, \ref{ch:heating}\\
ASTRA + SPIDER (Fenix) & перенесення + вільна межа + керування & 1.5D+2D & \cite{Fable2022} & \ref{ch:transport}, \ref{ch:wall}\\
TSC & вільна межа, кола котушок, перенесення & 1.5D+2D & \cite{Jardin2000} & \ref{ch:equilibrium}, \ref{ch:stability}\\
код домішок PPPL & зарядові стани $\times$ радіус & 1D$\times z$ & \cite{Hulse1992} & \ref{ch:transport}\\
Maxfea & GS з вільною межею & 2D & \cite{Merlo2011} & \ref{ch:equilibrium}\\
TASK/T2 & перенесення ядра й периферії в MSCS & 2D & \cite{Seto2014} & \ref{ch:geometry}, \ref{ch:flows}\\
SOLPS5.0 & B2.5 (рідина) + Eirene (нейтрали) & 2D & \cite{AhoMantila2011} & \ref{ch:wall}, \ref{ch:flows}\\
NOVA-W, PEST & лінійна ідеальна МГД & 2D & \cite{Ward1992,Hanada1990} & \ref{ch:stability}\\
JOREK & нелінійна МГД (ELM) + повні орбіти W & 3D & \cite{vanVugt2019} & \ref{ch:stability}, \ref{ch:transport}\\
GTC & гірокінетика, $\delta f$-PIC & 5D & \cite{Wang2006} & \ref{ch:transport}, \ref{ch:flows}\\
ASCOT, ERO, DUSTT & тестові частинки Монте-Карло на фоні плазми & 3D & \cite{AhoMantila2011,Pigarov2006} & \ref{ch:wall}\\
РК4/РК5, Боріс & повна орбіта частинки & 3D & \cite{Moskvitina2010,vanVugt2019} & \ref{ch:particles}\\
\bottomrule
\end{tabularx}
\end{table}

«Інтегроване» моделювання~--- це зшивання рівнів: у Fenix ASTRA, SPIDER і модель керування AUG рахуються разом~\cite{Fable2022}; у TASK/T2 цикл «перенесення в MSCS $\to$ струм $\Jtor(R,Z)$ $\to$ GS з вільною межею $\to$ нова MSCS» повторюють до збіжності, зберігаючи тороїдальний і полоїдальний потоки (схема FCT), тож $q(\rho)$ незмінне~\cite[с.~8--9]{Seto2014}.

\paragraph{Інтеграція «від геометрії».} Протилежний порядок показує власне ядро проєкту, створене для~JET~\cite{ProjJET}. Спершу складають базу розмірів машини, де кожне число має посилання на сторінку джерела (\texttt{machines/\allowbreak jet1975/\allowbreak dimensions.yaml}; стан 1983~р.~--- задокументовані зміни \texttt{machines/\allowbreak jet1983/\allowbreak deltas.yaml}). З неї геометричне ядро (\texttt{src/\allowbreak tokviz/\allowbreak machine}) будує тривимірні компоненти --- котушки, камеру, ярмо, лімітери, порти --- з автоматичними тестами на зазори й перетини. На ту саму геометрію «запікають» фізику: аналітичну рівновагу Соловйова, що мусить уміститися в камеру із заданими зазорами, і гофрування тороїдального поля з закону Біо--Савара по реальних котушках TF. Результат експортують у сцену Blender і веб-переглядач (glTF). Перевага такого порядку в тому, що суперечності джерел (наприклад, таблиця форми плазми, що виходить за стінку камери) виявляються вже під час збирання, а не на готовому рисунку. Ціна --- проста фізика (жодного перенесення чи еволюції), тож це інтеграція за геометрією, а не за фізикою (розд.~\ref{ch:jet}).

\section{Числові методи корпусу}\label{sec:11-methods}

\begin{itemize}[leftmargin=*]
\item \textbf{Неявна блоково-тридіагональна схема.} У коді домішок~\cite[с.~8]{Hulse1992} невідомі $n_z(r_j)$ зв'язані лише із сусідніми зонами (дивергенція) і сусідніми зарядовими станами (іонізація, рекомбінація). Матриця блоково-тридіагональна; неявний розв'язувач дає й стаціонар лінійної моделі, для неокласичних потоків додано предиктор--коректор. Неявна й TSC~\cite[с.~3]{Jardin2000}.
\item \textbf{МСЕ.} Maxfea: GS з вільною межею на трикутних елементах першого порядку~\cite[с.~31--32]{Merlo2011} (розд.~\ref{ch:equilibrium}). TASK/T2: рівняння в MSCS розв'язуються разом і неявно; повністю неявний МСЕ лише анонсовано~\cite[с.~8--9]{Seto2014}. Ітерація Пікара з функціями Гріна~--- розд.~\ref{ch:equilibrium}, \cite{ProjGSsolver}.
\item \textbf{Орбіти:} Боріс точно зберігає $|\vect v|$ у чисто магнітному полі, РК4/РК5~--- ні (розд.~\ref{ch:particles}, \cite{vanVugt2019,Moskvitina2010}).
\item \textbf{$\delta f$-PIC.} GTC розщеплює $f=f_0+\delta f$ і розв'язує рівняння для $\delta f$~\cite[с.~10]{Wang2006}. Сітка~--- вздовж силових ліній ($\kpar\ll k_\perp$)~\cite[с.~6]{Wang2006}; 4-точкове гіроусереднення точне до $k_\perp\rho_i\le2$~\cite[с.~9--10]{Wang2006}.
\item \textbf{Зіткнення Монте-Карло.} $D_{\mathrm{neo}}=\fsa{(\Delta r)^2}/2t$ з питч-кутового розсіяння ведучих центрів важких домішок~\cite[с.~6]{Wong1987}. ASCOT (оператори Фоккера--Планка) і ERO на фоні SOLPS~\cite[с.~52--53]{AhoMantila2011}; DUSTT (пил)~\cite[с.~5]{Pigarov2006}.
\item \textbf{Ермітова інтерполяція п'ятого порядку} (неперервні функція та дві похідні) прискорює обчислення швидкісних коефіцієнтів у ${\sim}30$ разів~\cite[с.~35]{Marchuk2025}.
\item \textbf{Обхід сингулярності ЗДР} у нижньогібридному резонансі: метод штрафу й аналітичне продовження в комплексну площину~\cite[с.~50--51]{Marchuk2025} (розд.~\ref{ch:heating}).
\item \textbf{Простір станів і зниження порядку.} Для LQG-керування формою RFX-mod лінійну модель розміру $n=250$ зведено ганкелевим методом до порядку $n_R=26$. Похибка систем $\mathcal T_0$, $\mathcal T_R$ обмежена хвостом ганкелевих сингулярних чисел: $\|\mathcal T_0-\mathcal T_R\|_\infty\le2\sum_{i>n_R}\sigma_i$~\cite[с.~97, 101]{Merlo2011}.
\end{itemize}

\section{Валідація проти експерименту}\label{sec:11-validation}

У корпусі «перевірено» означає чотири різні речі.
(1)~\emph{Часові сліди за заданим керуванням.} Fenix для AUG \#36440 відтворює керовані величини, але струми активної котушки й пасивного контуру до $t\approx\SI{0,45}{с}$ відхиляються сильно (автори припускають брак струмів на відкритих силових лініях)~\cite[с.~5]{Fable2022}. Коли $l_{i3}$ підігнано (зменшенням крайового бутстреп-струму), збігаються й струми котушок~\cite[с.~6--7]{Fable2022} (розд.~\ref{ch:transport}).
(2)~\emph{Магнітні сигнали за виміряними струмами котушок.} TSC для NSTX~\cite[с.~3--4]{Jardin2000}: ту саму згоду дають множник 1,5 до коефіцієнтів перенесення за $\Zeff=2{,}7$ і 1,0 за $\Zeff=5$~--- модель не визначено однозначно.
(3)~\emph{Локальні вимірювання.} SOLPS5.0 на AUG занижує паралельні іонні потоки в 2--3 рази в прямому полі~\cite[с.~61]{AhoMantila2011}, а зміну осадження $^{13}$C при оберненні поля відтворює лише якісно~\cite[с.~63]{AhoMantila2011} (розд.~\ref{ch:flows}, \ref{ch:wall}).
(4)~\emph{Код проти коду.} STREAM проти DYON на JET \#77210: для згоди в STREAM взято довжину зв'язку $L_f/3$, як у DYON~\cite[с.~13, 15]{Hoppe2022} (розд.~\ref{ch:disruptions}). Це верифікація, а не валідація.
Урок: незалежний тест~--- лише величина поза підгонкою.

\section{Реконструкція рівноваги за даними і машинне навчання}\label{sec:11-ml}

\textbf{Задача Fusion Equilibrium Challenge}~\cite{FEC2026}: відновити $\psi(R,Z)$ на сітці $65\times65$ за струмами котушок і томсонівським розсіянням, \emph{без} магнітних зондів; навчання~--- 7\,041 розряд DIII-D, MAST~--- без донавчання. Оцінка (оцінювач (scorer) \texttt{fusion\_scoring/common.py}):
\begin{equation}\label{eq:11-score}
\mathcal S=0{,}55\max(0,R^2_\psi)+0{,}15\max(0,R^2_{\qnf,\betaN})+0{,}10\,(1-D_{\mathrm{LCFS}})+0{,}20\,\mathrm{Consistency}.
\end{equation}
$D_{\mathrm{LCFS}}$~--- симетрична відстань Гаусдорфа між LCFS з поданого й істинного $\psi$, поділена на середній $R_{\mathrm{geo}}$ істинної LCFS і обрізана до 1 (безрозмірна); $R^2_{\qnf,\betaN}$~--- середнє двох об'єднаних (pooled) $R^2$; Consistency~--- середнє $\max(0,R^2_j)$ за сімома скалярами ($\Rax$, $\Zax$, $\kappa$, $\dup$, $\dlo$, об'єм, $\li$), які оцінювач обчислює з поданого $\psi$ і з істини (E27~\cite{ProjRegisters}). Усі 1\,206 розрядів MAST випуску (з FAIR-MAST, CC~BY-SA~4.0) видано під CC~BY~4.0; share-alike (§4(b)) цього, за нашим прочитанням, не дозволяє; до відповіді (Q13~\cite{ProjRegisters}) трактуємо їх як BY-SA.

\textbf{Базові моделі (baseline) проєкту.} У \texttt{experiments.py} (поділ за розрядами всередині 50 розрядів) MLP дає $R^2_\psi=0{,}915$, Ridge~--- $0{,}396$ (на 3 розрядах порядок інший: Ridge $0{,}963$, випадковий ліс $-2{,}49$). PCA+Ridge, навчений на 40 розрядах і оцінений на 8 \emph{інших}, дає $\mathcal S=0{,}176$ при $R^2_\psi=0{,}064$ і $1-D_{\mathrm{LCFS}}=0{,}90$ (власне відтворення, \texttt{SCORE.md}; оцінювач~\cite{FEC2026}). Числа не порівнянні: \texttt{experiments.py} усереднює $R^2$ по пікселях, оцінювач рахує об'єднаний $R^2$, а вибірки розрядів різні; 8 відкладених розрядів~--- мала вибірка. Гарна межа~--- не гарне $\psi$.

\textbf{Фізична нев'язка.} $R^2_\psi$ і Consistency міряють близькість до мітки EFIT, а не те, чи є $\psi$ рівновагою; нев'язку $g(\psi)$ (розд.~\ref{ch:equilibrium}, E6--E8) обчислюють із самого $\psi$.
\begin{established}{E22, E23}
E22 (прочитано): нейрооператори для GS мали $0{,}25\,\%$ відносної $L^2$-похибки за великих нев'язок (Ding та ін., arXiv:2511.19114). Навчене відображення не є проєкцією Гальоркіна, тож аналога леми Сеа немає (наш коментар). E23: жоден бенчмарк (FEC, TokaMark, PDEBench, The Well) не оцінює нев'язку PDE; Consistency~--- це $R^2$ похідних скалярів~\cite{ProjRegisters}.
\end{established}

\begin{refuted}{R11 (була гіпотеза C2)}
«Спектрально зважена втрата $\sum_k w(k)|\widehat{\delta\psi}(k)|^2$ з $w(k)$ за кривою чутливості розбору D1 покращує реальні моделі». Перевірено 2026-09-30 за критеріями, записаними до навчання: PCA+Ridge на вагу не реагує ($|\Delta|<10^{-4}$), у UNet\_Lite (3 зерна, 8 відкладених розрядів) усі ваги точково трохи гірші за пласку MSE, інтервали містять нуль~\cite{ProjRegisters,ProjGuide}. Побічно (E37): UNet\_Lite має $R^2_\psi=0{,}976$, але Consistency лише $0{,}179$~--- нижче, ніж PCA+Ridge з $R^2_\psi=0{,}090$ ($0{,}270$). Розчеплення, передбачене D1, видно на навченій моделі.
\end{refuted}

\begin{refuted}{R12 (половина гіпотези C4)}
«Додавання GS-члена ($g$, \texttt{GS\_score}) змінить порядок базових моделей». Перевірено 2026-10-01 на семи моделях \texttt{experiments.py} і PCA+Ridge: точковий порядок змінюється ($\tau_{\mathrm{Kend}}=0{,}857$), але жодна перестановка не тримається в 95\,\% бутстреп-повторень~\cite{ProjRegisters,ProjGuide}. Побічно (E38): мережі мають $R^2_\psi=0{,}96$--$0{,}98$, а нев'язку ГШ $g=0{,}81$--$0{,}96$~--- більшу, ніж в істини з 1\,\% шуму. Їхні карти потоку не є рівновагами, а гладкі лінійні прогнози проходять завдяки гладкості. GS-член природніше вживати як шлюз, а не як доданок.
\end{refuted}

\begin{conjecture}{C4, C6}
Нехай $\mathcal S'=\mathcal S+{}$GS-член${}+{}$член калібрування. \textbf{C4 (частина про калібрування):} додавання члена калібрування змінить порядок базових моделей. \textbf{Спростує:} порядок той самий. \textbf{Вартість:} пів дня після C6~\cite{ProjRegisters}. \textbf{Студенту:} знайти $\tau_{\mathrm{Kend}}$ між порядками моделей за $\mathcal S$ і за $\mathcal S'$.

\textbf{C6:} UQ-член~--- \emph{одночасні} конформні смуги на полі $65\times65$, а не 4\,225 поточкових~\cite{ProjRegisters}. Стандартний split-conformal: $\hat u$~--- масштаб похибки на піксель, підігнаний на окремій частині; для розряду $j$ $s_j=\max_{t,\mathrm{px}}|\psi-\hat\psi|/\hat u$; смуга $\hat\psi\pm\hat c\hat u$, де $\hat c$~--- порядкова статистика номер $\lceil(n_{\mathrm{cal}}+1)(1-\varepsilon_{\mathrm{cf}})\rceil$ з $n_{\mathrm{cal}}$ калібрувальних $s_j$; за обмінюваності розрядів покриття всього поля $\ge1-\varepsilon_{\mathrm{cf}}$. \textbf{Спростує:} покриття~--- лише за абсурдно широких смуг. \textbf{Вартість:} близько дня. 
\end{conjecture}

\begin{conjecture}{C9, C10 --- відкладено свідомо}
\textbf{C9:} перенесення DIII-D$\to$MAST покращить фізично коваріантне подання (коефіцієнти функцій Гріна замість пікселів); відкладено~--- це задача живого конкурсу. \textbf{C10:} частина похибки моделі~--- успадковане зміщення міток EFIT~\cite{ProjRegisters}. Спростування й вартість (наші): C9~--- $\mathcal S_{\mathrm{MAST}}/\mathcal S_{\mathrm{D3D}}$ не кращий, ніж з PCA; тижні. C10~--- спростування в межах корпусу неможливе (немає незалежної істини); $g$ дає лише нижню межу (шум дискретизації). \textbf{Студенту:} $g$(EFIT) на сітках $65\times65$ і $33\times33$~--- калібрування порогу $g$, а не перевірка C10.
\end{conjecture}

\begin{practicum}[базова модель і GS-нев'язка: \texttt{manual/code/ch11\_gs\_residual\_baseline.py}]
У \texttt{starter} через \texttt{.venv/bin/python}. \textbf{A} (мережа): \texttt{local\_score.py --mode perfect --n-shots 2} дає $\mathcal S=1$; далі \texttt{my\_experiments/baseline\_pca\_ridge.py --train --n-shots 40 --n-pca 50} і \texttt{local\_score.py --n-shots 8 --skip 60}; очікуйте $\mathcal S\approx0{,}18$. README згадує 10 розрядів і \texttt{--skip 7000}; числа~--- за \texttt{SCORE.md}. \textbf{B} (офлайн, 4~с; потребує артефакту з кроку A): \texttt{../../manual/code/ch11\_gs\_residual\_baseline.py}, 36 кадрів демо-розрядів. Медіани: $R^2_\psi\approx0{,}88$; $g(\text{EFIT})\approx0{,}010$, $g(\text{прогноз})\approx0{,}017$. Проєкція EFIT на ті самі 50 компонент: $1-R^2\approx2\cdot10^{-7}$, але $g\approx0{,}017$~--- нев'язку прогнозу задає базис PCA, а не регресія (порівн. E8: на 12 кадрах \#203702 PCA-5 дає 0,0093 проти 0,009 для EFIT --- інша вибірка й інший базис; чи є демо-розряди серед навчальних, невідомо). \textbf{C:} перенавчіть із \texttt{--n-pca} 5, 20, 100 (команда кроку A з \texttt{--n-pca N}; артефакт перезаписується) і перезапустіть B; побудуйте $g$ проти $1-R^2$.
\end{practicum}

\subsection*{Контрольні запитання}
\begin{enumerate}
\item Яке розділення масштабів виправдовує послідовність рівноваг GS? Коли воно хибне?
\item Чому згода TSC з NSTX за двох $(\text{множник},\Zeff)$ не валідує перенесення? Аналогія в EFIT?
\item Назвіть три причини, чому $R^2_\psi=0{,}40$ і $0{,}06$ для Ridge не суперечать одне одному.
\item Чому $g$ не бачить зміщення міток EFIT?
\end{enumerate}

\subsection*{Задачі}
\begin{enumerate}
\item \emph{Метрика.} PCA+Ridge (власне відтворення, \texttt{SCORE.md}): $R^2_\psi=0{,}0640$, $R^2_{\qnf,\betaN}=-0{,}7362$, $1-D_{\mathrm{LCFS}}=0{,}9035$. Додатні $R^2_j$ лише в $\Rax$ (0,4352), $\Zax$ (0,2715), $\kappa$ (0,4967), $\dup$ (0,5704); решта від'ємні. Знайдіть Consistency і $\mathcal S$ за~\eqref{eq:11-score}. Яке $R^2_\psi$ дало б $\mathcal S=0{,}5$ (інші члени незмінні)?
\item \emph{Вартість розв'язувачів.} (а) W: $N_Z=75$ зарядових станів (з нейтральним), $N_r=100$ зон (наші припущення; структура матриці~--- \cite[с.~8]{Hulse1992}): відношення вартості щільного LU $\tfrac23(N_rN_Z)^3$ до блокового Томаса $\approx5N_rN_Z^3$ (стандартні оцінки)? (б) Ріккаті $\propto n^3$: виграш від $n=250\to n_R=26$~\cite{Merlo2011}? Умова на хвіст $\sigma_i$?
\end{enumerate}
